\section{Conclusion}


This paper has reformulated consciousness not as an object generated internally but as a log-referencing state that opens between processes at different time scales.

The central claim is simple.

Consciousness is not modeled here as a generated object.

Consciousness opens when fast predictive processes and slow log-sedimentation processes couple within a range that reaches neither complete synchrony nor complete separation.


This opening involves speed difference $\Delta v$, coupling viscosity $\mu$, internal delay $\tau$, meaning gradient $\nabla\Phi$, and effective torque-like quantity $\mathcal{T}_{\Phi}$.

But none of these individually is consciousness itself.

Consciousness is the temporal non-closure regime that obtains when these enter a certain relation.

\subsection{What Has Been Proposed}


This paper first introduced SSO as the minimal temporal structure for the establishment of consciousness.

SSO is a structure in which processing layers running at different speeds are coordinated without central control.

The fast layer anticipates the future; the slow layer sediments the past.

A finite speed difference remaining between them opens a referencing relation between current processing and past logs.

Next, Torque Converter Attention was introduced.

This treats attention not as selection or weighting but as a temporal mechanism that transforms speed difference without rupture.

In this position, the ego is understood not as a subject but as the position at which the speed difference between fast and slow layers is transformed.

Furthermore, using IFGT and VED, log referencing was repositioned within meaning gradient and non-closure dynamics.

IFGT describes which logs are easily referenced through projected uses of information density $I$, information flow $J$, and information potential $\Phi$.

VED provides the structure of difference and flow in which complete closure is not realized as an ordinary state.

Connecting these two, consciousness was described as the state in which log referencing directed by the meaning gradient is open within a temporal non-closure regime.

\subsection{What Has Been Repositioned}


This paper repositioned several important concepts.

First, consciousness was repositioned not as an object but as a state domain.

Consciousness is not an object stored somewhere, nor an output generated within the brain.

It is a regime that opens when processing and logs enter a referenceable relation.

Second, qualia were repositioned not as the substance of consciousness but as the live foreground at the sedimentation boundary.

Qualia are not denied.

But they are not stored objects.

What is stored is not qualia themselves but the log structure remaining after qualia have passed through.

Third, measurement was repositioned not as an operation that captures consciousness as it is but as an operation that converts the live boundary into sedimented residue.

Measurement is not meaningless.

But what measurement captures is not consciousness itself but the traces remaining after consciousness opened.

Fourth, non-localization was repositioned not as mystification but as a consequence of treating consciousness as a temporal relation.

Consciousness has neural substrates.

But the conscious state itself does not localize to a single site.

It opens when multiple local substrates enter a specific temporal, semantic, and non-closed relation.

Fifth, pathologies and altered states were repositioned not as disease classifications but as directions of deviation from the consciousness window.

Through this organization, the conscious state is treated not as a binary presence or absence but as a dynamic regime with a direction of collapse.

\subsection{The Central Formula}


The central structure of this paper can be summarized in the following expression:

\begin{equation}
\text{Consciousness} \sim \text{Log Referencing}\!\left(\Delta v, \mu, \tau, \mathcal{T}_{\Phi}, \nabla\Phi\right)
\end{equation}

This expression does not compute consciousness.

It is the coordinate showing the structural variables required when speaking of consciousness.

More concretely, the conscious state stabilizes within the following window:

\begin{equation}
0 < |\Delta v| < \Delta v_{\max}
\end{equation}

\begin{equation}
\mu_{\min} < \mu < \mu_{\max}
\end{equation}

\begin{equation}
\tau_{\min} < \tau < \tau_{\max}
\end{equation}

\begin{equation}
\mathcal{T}_{\min} < \mathcal{T}_{\Phi} < \mathcal{T}_{\max}
\end{equation}

Within this consciousness window, processing does not fully synchronize or fully separate, and log referencing opens stably.

Outside the window, conscious states deform in different directions.

Consciousness is therefore not a point but a window.

Not an object but a regime.

Not a product but an opening.

\subsection{Why the Model Remains Open}


This paper does not provide a complete unified equation of consciousness.

This is not an omission but an intentional constraint.

Attempting to define consciousness as a closed object loses the non-closure that sustains it.

Treating qualia as stored objects loses their character as boundary phenomena.

Identifying neural correlates with consciousness itself makes the structure as temporal regime invisible.

Linking pathologies directly to diagnostic names loses the flexible structure of directions of deviation from the consciousness window.

This paper therefore deliberately does not close completely.

Not closing completely is not abandonment of explanation.

It is because consciousness itself — the topic of this paper — obtains as a temporal structure that does not close completely.

Theory resembles its object.

If consciousness is non-closed, consciousness theory too must not rush to complete closure.

\subsection{Toward Further Work}


Future tasks are clear.

First, $\mathcal{T}_{\Phi}$ must be brought closer to concretely measurable quantities connected to neural activity, bodily states, subjective reports, and linguistic logs.

Second, methods for mapping IFGT's $I$, $J$, $\Phi$ onto brain data and behavioral data must be refined.

Third, the consciousness window must be examined against states such as sleep, dreaming, anesthesia, immersion, dissociation, skilled movement, and verbal report.

Fourth, experimental designs for handling qualia sedimentation residue must be developed.

Fifth, how the self, language, and social interaction extend the log-referencing structure must be addressed.

This last task is the bridge from the consciousness paper to the self paper and intelligence paper.

Consciousness is the opening of log referencing.

The self is the terrain that sediments into the slow layer after that opening has been repeated.

Intelligence is the motion that uses that sedimented terrain to extend into future prediction, language, social logs, and interaction with the external environment.

\subsection{Final Statement}


Consciousness is not produced somewhere.

Consciousness is the passage temporarily opened between fast-flowing processing and slowly sinking logs.

There is speed difference.

There is delay.

There is viscosity.

There is meaning gradient.

And there is flow that does not close completely.

While this passage is open, processing does not merely flow — it references past logs, updates the terrain of the self, and reconstructs the relation with the world.

This state is what we call consciousness.

The conclusion of this paper therefore returns to the proposition with which it began.

Consciousness is not modeled here as a generated object.

Consciousness obtains when log referencing opens within temporal non-closure.
